% !TEX root = ../main.tex
%% Corrigés — Chapitre 11 — Vecteurs, droites et plans de l'espace
%% Tous les résultats ont été vérifiés par programme (calcul exact).

\corrigesChapitre{11}{Vecteurs, droites et plans de l'espace}

\begin{deuxcolonnes}

\begin{corrige}{1}{Opérations sur les vecteurs de l'espace}
\begin{questions}
  \item $\vec u+\vec v\,\res{(3\pv3\pv1)}$ ;\\ $3\vec u-2\vec v\,\res{(4\pv-11\pv13)}$ ;\\ $\vec u+\vec v+\vec w\,\res{(3\pv5\pv2)}$.
  \item Si $\vec v=k\vec u$, la première coordonnée donne $k=\frac12$ et la deuxième $k=-4$ : c'est impossible. Les vecteurs ne
    sont $\res{\text{pas colinéaires}}$.
  \item $\vect{AB}\,\res{(3\pv-3\pv2)}$ et $I\left(\frac{1+4}{2}\pv\frac{2-1}{2}\pv\frac{3+5}{2}\right)$, soit
    $\res{I\left(\tfrac52\pv\tfrac12\pv4\right)}$.
\end{questions}
\end{corrige}

\begin{corrige}{2}{Vecteurs directeurs et alignement}
\begin{questions}
  \item $\vect{AB}\,\res{(2\pv2\pv2)}$ et $\vect{AC}\,\res{(1\pv1\pv1)}$.
  \item $\vect{AB}=2\vect{AC}$ : les vecteurs sont colinéaires et les points $A$, $B$, $C$ sont $\res{\text{alignés}}$.
  \item $\vect{AB}\,(2\pv2\pv2)$ et $\vect{AC}\,(1\pv1\pv1)$ dirigent $(AB)$. Comme $\vect{AC}=\frac12\vect{AB}$, $C$ est le
    $\res{\text{milieu de }[AB]}$.
\end{questions}
\end{corrige}

\begin{corrige}{3}{Décompositions dans un cube}
\methode{relation de Chasles et égalités de vecteurs du cube : $\vect{BC}=\vect{AD}$, $\vect{CG}=\vect{AE}$…}
\begin{questions}
  \item $\vect{AG}=\vect{AB}+\vect{BC}+\vect{CG}$\\ $\phantom{\vect{AG}}=\res{\vect{AB}+\vect{AD}+\vect{AE}}$.\\
    $\vect{EC}=\vect{EA}+\vect{AC}=\res{\vect{AB}+\vect{AD}-\vect{AE}}$.\\
    $\vect{BH}=\vect{BA}+\vect{AD}+\vect{DH}$\\ $\phantom{\vect{BH}}=\res{-\vect{AB}+\vect{AD}+\vect{AE}}$.\\
    $\vect{DF}=\vect{DA}+\vect{AB}+\vect{BF}$\\ $\phantom{\vect{DF}}=\res{\vect{AB}-\vect{AD}+\vect{AE}}$.
  \item $\vect{FG}=\vect{AD}$ et $\vect{HD}=-\vect{AE}$, donc la somme vaut $\vect{AB}+\vect{AD}-\vect{AE}=\res{\vect{EC}}$.
  \item En ajoutant les trois expressions de la question 1, les coefficients de $\vect{AB}$, $\vect{AD}$, $\vect{AE}$ valent
    tous $1$ : la somme est $\vect{AB}+\vect{AD}+\vect{AE}=\vect{AG}$.
  \item $I$ est le milieu de $[BG]$ :\\ $\vect{AI}=\vect{AB}+\frac12\vect{BG}=\vect{AB}+\frac12\big(\vect{BC}+\vect{CG}\big)$,\\
    soit $\res{\vect{AI}=\vect{AB}+\tfrac12\vect{AD}+\tfrac12\vect{AE}}$.
\end{questions}
\end{corrige}

\begin{corrige}{4}{Configuration dans un cube}
\begin{questions}
  \item $A(0\pv0\pv0)$, $B(1\pv0\pv0)$, $C(1\pv1\pv0)$, $D(0\pv1\pv0)$, $E(0\pv0\pv1)$, $F(1\pv0\pv1)$, $G(1\pv1\pv1)$,
    $H(0\pv1\pv1)$.
  \item $\res{\vect{AG}=\vect{AB}+\vect{AD}+\vect{AE}}$, car $G(1\pv1\pv1)$.
  \item $\vect{AC}\,\res{(1\pv1\pv0)}$.
  \item $\vect{EF}\,(1\pv0\pv0)$ et $\vect{EH}\,(0\pv1\pv0)$ ne sont pas colinéaires et $E$, $F$, $H$ sont dans le plan $(EFG)$ :
    $\res{\text{oui}}$.
  \item $\vect{AC}\,(1\pv1\pv0)$, $\vect{AF}\,(1\pv0\pv1)$, $\vect{AH}\,(0\pv1\pv1)$. Si $\vect{AH}=a\vect{AC}+b\vect{AF}$, alors
    $a=1$ et $b=1$, mais $a+b=0$ : impossible. Les points ne sont $\res{\text{pas coplanaires}}$ ($ACFH$ est un tétraèdre
    régulier : ses arêtes sont des diagonales de faces).
\end{questions}
\end{corrige}

\begin{corrige}{5}{Combinaisons linéaires et coplanarité}
\begin{questions}
  \item Si $\vec v=k\vec u$, la première coordonnée donne $k=0$, donc $\vec v=\vec0$ : faux. $\vec u$ et $\vec v$ ne sont pas
    colinéaires : ils $\res{\text{dirigent un plan}}$.
  \item $\vec w=a\vec u+b\vec v$ : $a=2$ (abscisse), $b=1$ (cote) et l'ordonnée $2a+b=5$ est vérifiée. Donc
    $\res{\vec w=2\vec u+\vec v}$ : $\vec u$, $\vec v$, $\vec w$ sont coplanaires.
  \item $a=3$, $b=2$, mais $2a+b=8\neq7$ : $\vec t$ n'est $\res{\text{pas}}$ une combinaison de $\vec u$ et $\vec v$.
\end{questions}
\end{corrige}

\begin{corrige}{6}{Coplanarité de points}
\begin{questions}
  \item $\vect{AB}\,(2\pv1\pv0)$, $\vect{AC}\,(1\pv2\pv1)$, $\vect{AD}\,(3\pv3\pv1)$. La cote de $\vect{AB}$ est nulle mais pas
    celle de $\vect{AC}$ : ils ne sont pas colinéaires, et $A$, $B$, $C$ définissent un plan.
  \item $\vect{AD}=\vect{AB}+\vect{AC}$ : les points sont $\res{\text{coplanaires}}$, et $ABDC$ est un
    $\res{\text{parallélogramme}}$ (règle du parallélogramme).
  \item $\vect{AE}\,(1\pv0\pv2)=a\vect{AB}+b\vect{AC}$ : la cote donne $b=2$, l'abscisse $a=-\frac12$, mais l'ordonnée
    $a+2b=\frac72\neq0$. Donc $\res{E\notin(ABC)}$.
\end{questions}
\end{corrige}

\begin{corrige}{7}{Vecteurs coplanaires dans un cube}
\begin{questions}
  \item $\vect{IJ}=\vect{IF}+\vect{FB}+\vect{BJ}$\\ $\phantom{\vect{IJ}}=\frac12\vect{AB}-\vect{AE}+\frac12\vect{AD}$.\\
    Or $\vect{CE}=-\vect{AB}-\vect{AD}+\vect{AE}$ et $\vect{CG}=\vect{AE}$, donc\\
    $-\frac12\vect{CE}-\frac12\vect{CG}=\frac12\vect{AB}+\frac12\vect{AD}-\vect{AE}$ :\\
    c'est $\vect{IJ}$. Les trois vecteurs sont $\res{\text{coplanaires}}$.
  \item $I\left(\frac12\pv0\pv1\right)$, $J\left(1\pv\frac12\pv0\right)$ : $\vect{IJ}\left(\frac12\pv\frac12\pv-1\right)$,
    $\vect{CE}\,(-1\pv-1\pv1)$, $\vect{CG}\,(0\pv0\pv1)$, et l'on vérifie
    $-\frac12\vect{CE}-\frac12\vect{CG}=\left(\frac12\pv\frac12\pv-1\right)$.
\end{questions}
\end{corrige}

\begin{corrige}{8}{Vecteurs coplanaires dans un pavé droit}
\begin{questions}
  \item $C(a\pv b\pv0)$, $D(0\pv b\pv0)$, $E(0\pv0\pv c)$, $G(a\pv b\pv c)$ :\\ $\vect{CA}\,(-a\pv-b\pv0)$, $\vect{DE}\,(0\pv-b\pv c)$,
    $\vect{DG}\,(a\pv0\pv c)$, et $\vect{DE}-\vect{DG}=(-a\pv-b\pv0)$. Donc $\res{\vect{CA}=\vect{DE}-\vect{DG}}$ : coplanaires.\\
    Sans coordonnées : $\vect{DE}-\vect{DG}=\vect{GD}+\vect{DE}=\vect{GE}$, et $\vect{GE}=\vect{CA}$ car $ACGE$ est un
    rectangle.
  \item $I\left(a\pv0\pv\frac c2\right)$ : $\vect{AI}\left(a\pv0\pv\frac c2\right)$, $\vect{DF}\,(a\pv-b\pv c)$,
    $\vect{HE}\,(0\pv-b\pv0)$. Si $\vect{AI}=x\vect{DF}+y\vect{HE}$, l'abscisse donne $x=1$ et la cote $x=\frac12$ :
    $\res{\text{non coplanaires}}$.
\end{questions}
\end{corrige}

\begin{corrige}{9}{Un point hors d'un plan}
\begin{questions}
  \item $\vect{AB}\,(1\pv1\pv0)$, $\vect{AC}\,(-1\pv1\pv1)$, $\vect{AD}\,(2\pv2\pv1)$.
  \item Les cotes $0$ et $1$ empêchent la colinéarité : $A$, $B$, $C$ ne sont pas alignés et définissent le plan $(ABC)$.
  \item La cote donne $b=1$, l'abscisse $a-b=2$ soit $a=3$, mais alors l'ordonnée $a+b=4\neq2$ : $\res{\text{impossible}}$.
  \item $\vect{AB}$, $\vect{AC}$, $\vect{AD}$ ne sont pas coplanaires, donc $D\notin(ABC)$ : les quatre points ne sont pas
    coplanaires, et $ABCD$ est un $\res{\text{tétraèdre}}$.
\end{questions}
\end{corrige}

\begin{corrige}{10}{Positions relatives de droites et d'un plan}
\begin{questions}
  \item Les cotes de $\vec{v_1}$ et $\vec{v_2}$ ($0$ et $1$) empêchent la colinéarité : ils dirigent un plan. $\vec{u_1}=k\vec{u_2}$
    donnerait $k=2$ et $k=-1$ : $\vec{u_1}$, $\vec{u_2}$ non colinéaires.
  \item $\vec{u_1}=a\vec{v_1}+b\vec{v_2}$ : $a=2$, $b=-1$, et $a+b=1$ convient : $\res{\vec{u_1}=2\vec{v_1}-\vec{v_2}}$.
    Pour $\vec{u_2}$ : $a=1$, $b=2$, mais $a+b=3\neq-1$.
  \item $\vect{CA}\,(1\pv0\pv0)$ : $a=1$, $b=0$, mais $a+b=1\neq0$, donc $A\notin\mathcal{P}$. Comme $\vec{u_1}$ dirige
    $\mathcal{P}$, $d_1$ est $\res{\text{strictement parallèle}}$ à $\mathcal{P}$.
  \item $\vec{u_2}$ ne dirige pas $\mathcal{P}$ : $d_2$ est sécante à $\mathcal{P}$. Un point de $d_2$ s'écrit $(t\pv2-t\pv2t)$ et un
    point de $\mathcal{P}$, $(s\pv s+r\pv1+r)$. D'où $s=t$, $r=2t-1$ et $t+2t-1=2-t$, soit $t=\frac34$ :
    $\res{\left(\tfrac34\pv\tfrac54\pv\tfrac32\right)}$.
  \item Un point de $d_1$ s'écrit $(1+2k\pv k\pv1-k)$. Les deux premières équations donnent $k=\frac13$ et $t=\frac53$, mais
    alors $1-k=\frac23\neq2t$. Pas de point commun ni de parallélisme : $d_1$ et $d_2$ sont $\res{\text{non coplanaires}}$.
\end{questions}
\end{corrige}

\begin{corrige}{11}{Section d'un tétraèdre par un plan de milieux}
\begin{questions}
  \item Dans le triangle $ABC$, $I$ et $J$ sont les milieux de $[AB]$ et $[AC]$ : d'après le théorème des milieux, $(IJ)\parallel(BC)$
    et $IJ=\frac12BC$. De même, dans le triangle $BCD$, $(LK)\parallel(BC)$ et $LK=\frac12BC$.
  \item Dans le triangle $ACD$, $(JK)\parallel(AD)$ et $JK=\frac12AD$ ; dans le triangle $ABD$, $(IL)\parallel(AD)$ et
    $IL=\frac12AD$.
  \item Les droites $(IJ)$ et $(LK)$ sont parallèles, donc coplanaires : $I$, $J$, $K$, $L$ sont $\res{\text{coplanaires}}$. De plus
    $\vect{IJ}=\frac12\vect{BC}=\vect{LK}$ : $IJKL$ est un $\res{\text{parallélogramme}}$.
  \item La section est le parallélogramme $IJKL$. Si $AD=BC$, alors $IJ=JK$ : deux côtés consécutifs sont égaux, la section est un
    $\res{\text{losange}}$.
\end{questions}
\end{corrige}

\begin{corrige}{12}{Vecteurs coplanaires dans un tétraèdre}
\begin{questions}
  \item Dans le triangle $ABC$, $\vect{IJ}=\frac12\vect{BC}$, d'où $\vect{BC}=2\vect{IJ}=2\vect{IJ}+0\,\vect{JK}$ : les trois
    vecteurs sont $\res{\text{coplanaires}}$.
  \item On a $B(1\pv0\pv0)$, $C(0\pv1\pv0)$ et $D(0\pv0\pv1)$, puis $I\left(\frac12\pv0\pv0\right)$, $J\left(0\pv\frac12\pv0\right)$ et
    $K\left(0\pv0\pv\frac12\right)$. Donc $\vect{BI}\left(-\frac12\pv0\pv0\right)$, $\vect{CJ}\left(0\pv-\frac12\pv0\right)$ et
    $\vect{DK}\left(0\pv0\pv-\frac12\right)$.
  \item Si $\vect{DK}=a\vect{BI}+b\vect{CJ}$, la cote donne $-\frac12=0$ : impossible. Les vecteurs ne sont $\res{\text{pas
    coplanaires}}$. En effet, ils sont colinéaires à $\vect{AB}$, $\vect{AC}$ et $\vect{AD}$, qui ne sont pas coplanaires puisque
    $ABCD$ est un tétraèdre.
\end{questions}
\end{corrige}

\end{deuxcolonnes}

%% ---------------------------------------------------------------------
\corrigesSouples
\begin{corrige}{13}{Le théorème du toit dans un cube}
\begin{questions}
  \item $\vect{BC}=\vect{EH}$ (ces deux vecteurs sont égaux à $\vect{AD}$) : $BCHE$ est un parallélogramme, ses sommets sont
    coplanaires. De même $\vect{AD}=\vect{FG}$ : $ADGF$ est un parallélogramme.
  \item $K$, centre du carré $ABFE$, est le milieu de ses diagonales $[BE]$ et $[AF]$. Ainsi $K\in(BE)\subset(BCE)$ et
    $K\in(AF)\subset(ADG)$.
  \item Le point $C$ n'est pas dans le plan $(ADG)$, qui est le plan $(ADGF)$ : les deux plans sont distincts. Ils ont le point $K$ en
    commun : ils sont sécants suivant une droite $\Delta$ passant par $K$. Le plan $(BCE)$ contient $(BC)$, le plan $(ADG)$ contient
    $(AD)$, et $(BC)\parallel(AD)$ : d'après le théorème du toit, $\Delta$ est $\res{\text{la parallèle à }(BC)\text{ passant par }K}$.
  \item $L$, centre du carré $DCGH$, est le milieu de $[CH]$ et de $[DG]$ : $L\in(CH)\subset(BCE)$ et $L\in(DG)\subset(ADG)$. Donc
    $L\in\Delta$, et $\res{\Delta=(KL)}$.
  \item $K\left(\frac12\pv0\pv\frac12\right)$ et $L\left(\frac12\pv1\pv\frac12\right)$ : $\vect{KL}\,(0\pv1\pv0)=\vect{BC}$. De plus
    $\vect{BK}=\frac12\vect{BE}$, $\vect{CL}=\frac12\vect{CH}$, $\vect{AK}=\frac12\vect{AF}$ et $\vect{DL}=\frac12\vect{DG}$, ce qui
    confirme les appartenances.
\end{questions}
\end{corrige}

\begin{corrige}{14}{Section d'un cube par un plan}
\begin{codecote}[0.63]
\begin{questions}
  \item Dans la face $ADHE$, $(IJ)$ joint les milieux de $[AD]$ et $[AE]$ : elle est parallèle à $(DE)$ et coupe $(EH)$ en un point
    $P$ situé hors du cube, au-delà de $E$, avec $EP=\frac12EH$.
  \item $P\in(IJ)\subset(IJK)$ et $P\in(EH)\subset(EFG)$ ; de même, $K\in[FG]$ appartient aux deux plans, qui sont distincts. Leur
    intersection est $\res{\text{la droite }(PK)}$ ; elle coupe $[EF]$ en son milieu $R$.
  \item On relie $J$ et $R$ (face $ABFE$), puis $R$ et $K$ (face $EFGH$). Dans la face $BCGF$, parallèle à $ADHE$, la parallèle à
    $(IJ)$ passant par $K$ coupe $[CG]$ en son milieu $N$. Dans la face $ABCD$, parallèle à $EFGH$, la parallèle à $(RK)$ passant par
    $I$ coupe $[CD]$ en son milieu $M$. La section est l'hexagone $\res{IJRKNM}$.
  \item $\vect{IJ}\left(0\pv-\frac12\pv\frac12\right)$, $\vect{IK}\,(1\pv0\pv1)$, $\vect{IM}\left(\frac12\pv\frac12\pv0\right)$. Si
    $\vect{IM}=a\vect{IJ}+b\vect{IK}$, l'abscisse donne $b=\frac12$, l'ordonnée $a=-1$, et la cote $\frac a2+b=0$ est vérifiée :
    $M\in(IJK)$. Chaque côté joint les milieux de deux arêtes voisines d'une face et mesure $\frac{\sqrt2}{2}$ : la section est
    un $\res{\text{hexagone régulier}}$.
\end{questions}
\cote
\centering
\begin{tikzpicture}[otEspace,scale=2.15]
  \otPaveCoord{1}{1}{1}
  \coordinate (I) at (0,0.5,0); \coordinate (J) at (0,0,0.5); \coordinate (K) at (1,0.5,1);
  \coordinate (R) at (0.5,0,1); \coordinate (N) at (1,1,0.5); \coordinate (M) at (0.5,1,0);
  \coordinate (P) at (0,-0.5,1);
  \fill[otGreen,opacity=0.22] (I) -- (J) -- (R) -- (K) -- (N) -- (M) -- cycle;
  \otPaveAretes\otPaveNoms
  \draw[otConstr] (I) -- (P) -- (E);
  \draw[otConstr] (P) -- (R);
  \draw[otGreenDark,line width=0.9pt] (J) -- (R) -- (K) -- (N);
  \draw[otGreenDark,line width=0.9pt,dash pattern=on 2.2pt off 1.6pt] (N) -- (M) -- (I) -- (J);
  \foreach \p/\a in {I/south east,J/east,K/south east,R/north west,N/west,M/south}{\node[otPtV] at (\p) {};
    \node[otSommet,anchor=\a] at (\p) {$\p$};}
  \node[otPt] at (P) {}; \node[otSommet,anchor=east] at (P) {$P$};
\end{tikzpicture}\\[0.5mm]
\legende{En orange, les traits de construction.}
\end{codecote}
\end{corrige}
